\section{Positive matrix realizations in every dimension}\label{sec:matrix57}
Let $d\ge2$ and let $\mathcal H_d^0$ be the real Euclidean space of traceless Hermitian $d\times d$ matrices, with scalar product $\langle A,B\rangle=\operatorname{tr}(AB)$ and Frobenius norm $\|\cdot\|_{\rm F}$. Let
\[
 \mathcal D_d=\{D=D^*\succeq0:\operatorname{tr}D=1\},\quad
 c_d=I/d,\quad r_d=\sqrt{1-1/d}.
\]
This section concerns numerical outputs in $\mathcal D_d$. The requirement that every decoder value be positive with trace one is part of the model; it is not imposed merely on the final average.

\begin{lemma}[Small balls on the complex Grassmannian]\label{lem:grassmann57}
Let $P$ be a fixed rank-$m$ orthogonal projection in $\C^d$, with $1\le m<d$. For Haar $U\in\mathrm{SU}(d)$ there is $C_{d,m}<\infty$ such that
\[
 \Pr\{\|UPU^*-Q\|_{\rm F}<h\}\le C_{d,m}h^{2m(d-m)}\quad(0<h\le1)
\]
for every rank-$m$ projection $Q$. For rank one, writing
$\operatorname{dist}_{\rm pr}(P,Q)=\sqrt{1-\operatorname{tr}(PQ)}$, one has $\operatorname{dist}_{\rm pr}(P,Q)=\norm{P-Q}_{\rm F}/\sqrt2$ and exactly
\begin{equation}\label{eq:projectivecap57}
 \Pr\{\operatorname{dist}_{\rm pr}(UPU^*,Q)<h\}=h^{2(d-1)}
                         \qquad(0\le h\le1).
\end{equation}
\end{lemma}
\begin{proof}
Conjugation is transitive, so fix $Q=\operatorname{diag}(I_m,0)$. A projection sufficiently close to $Q$ has range the graph of a unique matrix $Z\in\C^{(d-m)\times m}$, with
\[
 P(Z)=\begin{pmatrix}I\\Z\end{pmatrix}(I+Z^*Z)^{-1}(I\quad Z^*).
\]
This notation means the product of the three displayed matrices. The map and its inverse are smooth near zero, and its derivative at zero is $Z\mapsto\left(\begin{smallmatrix}0&Z^*\\Z&0\end{smallmatrix}\right)$, of Frobenius norm $\sqrt2\|Z\|_{\rm F}$. The invariant probability measure is normalized Riemannian volume on this compact homogeneous manifold; in this chart its density is smooth and bounded on a sufficiently small closed neighborhood. The inverse chart puts a Frobenius $h$-ball inside a Euclidean $Ch$-ball in real dimension $2m(d-m)$. This proves the bound for small $h$; increasing the constant covers the remaining interval. Transitivity makes the constant independent of $Q$.

For rank one choose $Q=e_1e_1^*$. A uniform unit vector can be obtained by normalizing independent standard complex Gaussian coordinates. If $X$ is the squared modulus of the first coordinate before normalization and $Y$ is the sum of the other squared moduli, then $X$ is exponential of mean one and $Y$ is an independent sum of $d-1$ such variables. For $0<z<1$,
\[
 \Pr\left\{\frac{X}{X+Y}>1-z\right\}
 =\E\exp\left(-\frac{1-z}{z}Y\right)=z^{d-1}.
\]
Put $z=h^2$. The endpoint cases follow by continuity. This proves \eqref{eq:projectivecap57}, and also shows that $\operatorname{dist}_{\rm pr}$ is the Frobenius chordal distance divided by $\sqrt2$.
\end{proof}

\begin{lemma}[A uniform conjugacy-orbit exponent]\label{lem:conjugacysmall57}
There is $C_d<\infty$ such that
\begin{equation}\label{eq:conjugacycap57}
 \sup_{\substack{A,B\in\mathcal H_d^0\\\|A\|_{\rm F}=\|B\|_{\rm F}=1}}
 m_G\{U:\|UAU^*-B\|_{\rm F}<h\}
                 \le C_d h^{2(d-1)}\qquad(0<h\le1).
\end{equation}
The exponent cannot be increased uniformly over $A$ and $B$.
\end{lemma}
\begin{proof}
Write the eigenvalues of $A$ as $a_1\ge\cdots\ge a_d$. Since their sum is zero and their squared sum is one, $a_1-a_d\ge d^{-1/2}$. Some adjacent gap therefore satisfies
\begin{equation}\label{eq:separatedcut57}
 a_m-a_{m+1}\ge\gamma_d:=\frac1{\sqrt d(d-1)}.
\end{equation}
Let $P_A$ be the spectral projection onto the first $m$ eigenvectors. It is well defined despite possible multiplicities within either block.

Two conjugates $A',A''$ of $A$, with eigenbases $x_i,y_j$ in this order, obey the exact identity
\begin{equation}\label{eq:spectralidentity57}
 \|A'-A''\|_{\rm F}^2
       =\sum_{i,j}(a_i-a_j)^2|\langle x_i,y_j\rangle|^2.
\end{equation}
Expanding both traces proves it. The cross-block terms and \eqref{eq:separatedcut57} imply
\begin{equation}\label{eq:projectorbound57}
 \|P_{A'}-P_{A''}\|_{\rm F}
                    \le\gamma_d^{-1}\|A'-A''\|_{\rm F},
\end{equation}
since the squared norm on the left is exactly the sum of the two cross-block overlap sums.

If the set in \eqref{eq:conjugacycap57} is nonempty, choose one member $U_0$. Every other member has $\|UAU^*-U_0AU_0^*\|_{\rm F}<2h$. Thus its rank-$m$ spectral projection lies within $2h/\gamma_d$ of $U_0P_AU_0^*$. The Haar law of $UP_AU^*$ is the invariant Grassmannian law. Lemma~\ref{lem:grassmann57} bounds the measure by
$C_{d,m}(2h/\gamma_d)^{2m(d-m)}$ for $h\le\gamma_d/2$. Since $m(d-m)\ge d-1$, and there are finitely many $m$, this gives \eqref{eq:conjugacycap57} with a constant independent of all eigenvalue multiplicities. Larger $h$ are covered by increasing $C_d$.

For sharpness take $A=B=(P-c_d)/r_d$ with $P$ rank one. By \eqref{eq:projectivecap57}, for sufficiently small $h$ the probability is
\[
                    (r_d^2h^2/2)^{d-1}.
\]
It cannot be bounded by a constant times $h^{p'}$ as $h\downarrow0$ for any $p'>2(d-1)$.
\end{proof}
The spectral split, rather than an assumption that all centroids remain rank one, supplies the uniformity needed by the stochastic contraction. The conjugation representation has ambient dimension $d^2-1$, but its worst small-ball exponent is $2(d-1)$.

\begin{theorem}[Unitary matrix-orbit width]\label{thm:matrixwidth57}
Let a finite alphabet in $\mathrm{SU}(d)$ contain the identity and support a probability law satisfying \eqref{eq:groupgap57}. Fix a rank-one projection $P_0$. For $0<a\le1$, set
\begin{equation}\label{eq:matrixexperiment57}
 F_a(U)=c_d+a(UP_0U^*-c_d),\qquad Y=\mathcal D_d,
\end{equation}
with Frobenius error. Its exact bounded-width threshold is $a r_d$, attained by the one-label decoder $c_d$. For $0\le\varepsilon<a r_d$, put $\zeta=a-\varepsilon/r_d$. Every realization satisfies
\begin{equation}\label{eq:matrixoccupation57}
 \#\{1\le t\le N:|S_t|\le k\}
 \le C\log(k+2)\left(1+k^{1/(d-1)}\log\frac1\zeta\right).
\end{equation}
For fixed $0<a<1$ and $0\le\varepsilon<a r_d$,
\begin{equation}\label{eq:matrixrate57}
 c\left(\frac{N}{\log(N+2)}\right)^{d-1}
            \le W_{N,\varepsilon}\le C(N+1)^{d-1},
\end{equation}
and the upper realization can be exact. The same orders hold for $a=1$ and every fixed $0<\varepsilon<r_d$.

More explicitly, for $a<1$ and $\varepsilon=0$ set $a'=a$; for $\varepsilon>0$ set $a'=a-\varepsilon/(2r_d)$. Whenever $a'<1$, an error at most $\varepsilon/2$ realization has
\begin{equation}\label{eq:matrixupper57}
 K\le C_d\left(1+\frac{N+1}{-\log a'}\right)^{d-1}
\end{equation}
labels. In the zero-error case this realization is exact. In \eqref{eq:matrixoccupation57} the constant depends only on $d$ and the spectral-gap data, not on the amplitude or tolerance. Extra finitely many unitary commands are allowed in both bounds.
\end{theorem}
\begin{proof}
Center at $c_d$ and divide all outputs by $r_d$. The resulting legal set is
\[
 Q=\operatorname{conv}\{(P-c_d)/r_d:P\text{ rank one}\}\subseteq\mathbb B_{\mathcal H_d^0}.
\]
The last inclusion follows from $\operatorname{tr}D^2\le1$ for $D\in\mathcal D_d$. The orbit spans $\mathcal H_d^0$: if a traceless Hermitian $A$ has $\operatorname{tr}(AP)=0$ for every rank-one $P$, then $v^*Av=0$ for every $v$, hence $A=0$. Conjugation-fixed matrices are scalar, so there are no fixed vectors in this traceless space. The orbit is complex projective space, of real dimension $2(d-1)$. Lemma~\ref{lem:conjugacysmall57} and Theorem~\ref{thm:uniformorbit57}, with $p=q=2(d-1)$, prove the lower and occupation bounds, as well as the radius statement after multiplying the norm by $r_d$.

We give a direct inner approximation to make positivity of the upper rows transparent. Choose a maximal $\delta$-separated family of rank-one projections $P_i$ for $\operatorname{dist}_{\rm pr}$, containing $P_0$. Its $\delta/2$ balls are disjoint, so \eqref{eq:projectivecap57} gives
\[
                         K\le(2/\delta)^{2(d-1)}.
\]
For any traceless Hermitian $A$, put $M=\lambda_{\max}(A)$. When $A\ne0$, $M>0$ and $\lambda_{\min}(A)\ge-(d-1)M$. Approximate a top eigenline by $P_i$ at distance at most $\delta$. Spectral decomposition then gives
\begin{equation}\label{eq:matrixsupport57}
 \operatorname{tr}(AP_i)\ge M-(M-\lambda_{\min}(A))\delta^2
                         \ge(1-d\delta^2)M.
\end{equation}
The trivial $A=0$ case is automatic. Comparing support functions in $\mathcal H_d^0$, with $K_0=\mathcal D_d-c_d$, yields
\[
 (1-d\delta^2)K_0\subseteq\operatorname{conv}\{P_i-c_d\}\subseteq K_0.
\]
Put $r=1-d\delta^2>0$. For each actual command $U_b$ choose one row of barycentric coordinates satisfying
\begin{equation}\label{eq:matrixrow57}
 \sum_jT_b(i,j)(P_j-c_d)=rU_b(P_i-c_d)U_b^*.
\end{equation}
Rows are nonnegative and sum to one by construction. Initialize at $P_0$ and decode label $i$ as
\[
                  D_i=c_d+(a'/r^N)(P_i-c_d).
\]
This is in $\mathcal D_d$ when $r^N\ge a'$: it is a convex combination of $c_d$ and $P_i$. Choosing
\[
 \delta^2=\min\left(\frac1{4d},\frac{-\log a'}{2d(N+1)}\right)
\]
ensures this inequality. Induction in \eqref{eq:matrixrow57} gives mean output exactly $F_{a'}(u_w)$. The target discrepancy is $|a-a'|r_d=\varepsilon/2$ when $\varepsilon>0$. The cardinal bound proves \eqref{eq:matrixupper57}. The case $a=1,\varepsilon=0$ is deliberately excluded from this dilation construction and is treated exactly in the next section.
\end{proof}
The family contains the qubit sphere as $d=2$, but for $d\ge3$ the pure-state orbit is a proper curved subset of the unit sphere in $\mathcal H_d^0$. Thus \eqref{eq:matrixrate57} does not follow by substituting the ambient dimension into the spherical theorem. Its logarithmic gap is explicit; no sharp constant or disappearance of that gap is claimed.

\begin{corollary}[Algebraic generating alphabets]\label{cor:algebraicunitary57}
Theorem~\ref{thm:matrixwidth57} applies to the identity together with any finite symmetric algebraic-entry set generating a dense subgroup of $\mathrm{SU}(d)$, equipped with a lazy symmetric law. Such a finite alphabet exists for every $d\ge2$.
\end{corollary}
\begin{proof}
The spectral gap is the imported theorem of Bourgain--Gamburd \cite{BG2012}. Laziness makes its one-sided gap an absolute $L^2$ contraction on mean-zero functions. Existence can be seen by placing two algebraic irrational-angle $\mathrm{SU}(2)$ rotations in every adjacent pair of coordinate directions. For example, use the quaternion lifts in the explicit adjacent-block construction below within each pair. Each pair generates its full $\mathrm{SU}(2)$ subgroup. Their Lie algebras contain the adjacent off-diagonal matrix units and their brackets, which generate $\mathfrak{su}(d)$. The connected compact closure is consequently $\mathrm{SU}(d)$. All entries remain algebraic. The spectral-gap constant is not computed by the finite regression checks.
\end{proof}

\subsection{A finite whole-law encoding}
Let $q_0=d^2-1$ and choose a Frobenius-orthonormal basis $H_1,\ldots,H_{q_0}$ of $\mathcal H_d^0$. Since $\|H_l\|_{\rm op}\le1$, the $2q_0$ matrices
\[
 E_{l,\pm}=\frac{I\pm H_l/2}{2q_0}
\]
are positive and sum to $I$. Put $\Phi(D)_{l,\pm}=\operatorname{tr}(E_{l,\pm}D)$ and $Y_\Phi=\Phi(\mathcal D_d)$. Every entry is at least $1/(4q_0)$. The map is injective, because
\begin{equation}\label{eq:tomography57}
 D-c_d=\sum_l2q_0\bigl(\Phi(D)_{l,+}-\Phi(D)_{l,-}\bigr)H_l.
\end{equation}
On the vector space parallel to the affine span of $Y_\Phi$, use the norm of the reconstructed matrix on the right. Then $\Phi$ is an affine isometry, and replacing each legal matrix decoder by its image preserves every clocked or stationary width exactly. This is one whole categorical law with a compact convex legality constraint, not separate binary event tolerances. For ordinary total variation the explicit norm comparison is
\[
 \frac{\|D-D'\|_{\rm F}}{4q_0}
 \le\TV(\Phi(D),\Phi(D'))
 \le\frac{\|D-D'\|_{\rm F}}{4\sqrt{q_0}}.
\]
It follows directly by comparing the $\ell^1$ and $\ell^2$ norms of the basis coefficients. Enlarging the legal output set from $Y_\Phi$ to the entire simplex is a different realization problem. In particular the extreme-output argument below cannot be transferred to that enlarged set: its target vectors are interior points of the simplex.

\begin{remark}[An explicit algebraic alphabet in every dimension]
In each adjacent coordinate plane $(j,j+1)$ embed the two matrices
\[
 \frac15\begin{pmatrix}-3+4i&0\\0&-3-4i\end{pmatrix},\qquad
 \frac15\begin{pmatrix}-3&4\\-4&-3\end{pmatrix},
 \qquad 1\le j<d,
\]
and act identically on all other coordinates. Include their inverses and
the identity. Each pair has infinite-order eigenvalues, since its trace
$-6/5$ is not an integer, and its two torus closures generate the adjacent
$SU(2)$. The adjacent skew-Hermitian matrix units and their brackets
span $\mathfrak{su}(d)$. The closed generated subgroup is therefore
$SU(d)$. Bourgain--Gamburd is applied to this actual finite algebraic
set, not merely to the existence of an unspecified set. The resulting
spectral-gap constant is not numerically evaluated here.
\end{remark}
The Frobenius chordal distance between rank-one projections is $\sqrt2$
times the projective distance used in the net construction. In the
conjugacy-cap comparison, choosing one member of a nonempty target ball
avoids imposing any spectral condition on its arbitrary center. All
adjacent-gap and cap constants here depend on $d$. The amplitude
$ a'=a-\varepsilon/(2r_d)$ in the positive-error upper bound deliberately
spends only half the available error.
